\documentclass{article}

% -------------------------------------------------
% Packages
% -------------------------------------------------

\usepackage{xcolor}
\usepackage{amsmath}
\usepackage{amssymb}
\usepackage{amsthm}
\usepackage{graphicx}
\usepackage{subcaption}
\usepackage{matlab-prettifier}
\usepackage{listings}
\usepackage{blindtext}
\usepackage{geometry}
\usepackage{float}
\usepackage{dirtytalk}
\usepackage{enumitem}
\usepackage{derivative}
\usepackage{biblatex}
\usepackage{svg}
\usepackage{hyperref}

% -------------------------------------------------
% Colors
% -------------------------------------------------

\definecolor{darkgreen}{rgb}{0.0, 0.5, 0.0}

% -------------------------------------------------
% Hyperlinks
% -------------------------------------------------

\hypersetup{
    colorlinks=true,
    linkcolor=blue,
    citecolor=blue,
    urlcolor=blue
}

% -------------------------------------------------
% Theorem environments
% -------------------------------------------------

% Numbered
\newtheorem{theorem}{Theorem}[section]
\newtheorem{lemma}[theorem]{Lemma}
\newtheorem{proposition}[theorem]{Proposition}

% Unnumbered
\newtheorem*{theorem*}{Theorem}
\newtheorem*{lemma*}{Lemma}
\newtheorem*{proposition*}{Proposition}

\theoremstyle{definition}
\newtheorem{definition}[theorem]{Definition}
\newtheorem{example}[theorem]{Example}

\theoremstyle{remark}
\newtheorem{remark}[theorem]{Remark}

% -------------------------------------------------
% Equation numbering
% -------------------------------------------------

% Reset equation numbering at each subsection
\makeatletter
\@addtoreset{equation}{subsection}
\makeatother

% If you want equations to be numbered by subsection,
% uncomment the following line:
% \renewcommand{\theequation}{\thesubsection.\arabic{equation}}

% -------------------------------------------------
% Python / MATLAB code formatting
% -------------------------------------------------

\lstset{
    language=Python,
    basicstyle=\ttfamily\small,
    keywordstyle=\color{blue},
    stringstyle=\color{darkgreen},
    commentstyle=\color{gray},
    showstringspaces=false,
    frame=single,
    numbers=left,
    numberstyle=\tiny\color{gray},
    breaklines=true
}

% -------------------------------------------------
% Page geometry
% -------------------------------------------------

\geometry{
    a4paper,
    total={160mm,250mm},
    left=25mm,
    top=30mm,
    footskip=10mm
}

% -------------------------------------------------
% Document settings
% -------------------------------------------------

\setlength{\parindent}{0pt}

% -------------------------------------------------
% Title
% -------------------------------------------------

\title{Monte Carlo Simulation Uge 1}
\author{Thomas Sejer Canham}
\date{\today}

% =================================================
% Document
% =================================================

\begin{document}

In week 1 we discussed some behind the scenes of Monte Carlo Simulation. 
\section{Linear Congruential Generator}

An important part of Monte Carlo Simulation is the generation of pseudo-random numbers. One of the first techniques for generating pseudo-random numbers was the Linear Congruential Generator (LCG). The LCG is defined by the recurrence relation:
\begin{equation}
    Z_{n+1} = (aZ_n + c) \mod m
\end{equation}
Where specific values of $a$, $c$, and $m$ are chosen to produce a sequence of pseudo-random numbers and it is initialized with some $Z_0 \in \mathbb{N}$. The choice of these parameters is crucial for the quality of the generated numbers. If chosen poorrly, the cycle length can be very short, on the order of 1 or 2, which is not desirable. However, number theory gives results that improve this to the order of $m$. The number theory results are beyond the scope of this course.

\section{Monte Carlo Integration in 1D}
We can use Monte Carlo Simulation to estimate the value of an integral. Let $h(x)$ be defined on the unit interval $(0,1)$. We can estimate the integral of $h(x)$ by generation $N$ random numbers $x_1, x_2, \ldots, x_N$ uniformly distributed in $(0,1)$ and computing the average of $h(x_i)$:
\begin{equation}
   \hat{\ell} = \int_0^1 h(x) dx \approx \frac{1}{N} \sum_{i=1}^{N} h(x_i)
\end{equation}
As the law of large numbers ensures convergence of the average to the expected value, we have that
\begin{equation}
    \dfrac{1}{N}\sum_{i=1}^{N} h(x_i) \to \mathbb{E}[h(X_i)] = \int_0^1 h(x)f(x) dx \text{ for } N \to \infty.
\end{equation}
Since the randoms numbers i i.i.d and uniformly distributed, we have that $f(x) = 1$ for $x \in (0,1)$ and $0$ otherwise. Thus we have that
\begin{equation}
    \dfrac{1}{N}\sum_{i=1}^{N} h(x_i) \to \int_0^1 h(x) dx \text{ for } N \to \infty.
\end{equation}

\section{Testing for Statistical Properties of Pseudo Random Numbers (PRNs)}
When we have indeed made a PRN Generator, we want to test it for statistical properties. We can do this by generating a large number of PRNs and then testing them for uniformity and independence. To test uniformity, we discretize the unit interval into $r$ bins, and count the number of points in each bin. We compare this to the expected number of points. Practially this is done using the $\chi^2$ distributed test statistic:
\begin{equation}
    T_N(u_1, \ldots, u_N) = \sum_{i=1}^{r} \frac{(N_i - N/r)^2}{N/r}
\end{equation}
Where $N_i$ is the number of points in bin $i$. If the PRNs are indeed uniformly distributed, then $T_N$ is $\chi^2$ distributed with $r-1$ degrees of freedom. We can then compare the value of $T_N$ to the critical value of the $\chi^2$ distribution to determine if we reject or accept the null hypothesis that the PRNs are uniformly distributed at a given significance level. See section 1.4 of the lecture notes.

Somewhat more involved is the test for autocorrelations by the lengths of increasing subsequences in the PRN sequence. See example 1.6 of the lecture notes.D

\section{Output Analysis}


\end{document}